\documentclass[spanish,notitlepage,letterpaper, 12pt]{article} 
\usepackage[spanish]{babel} 
\usepackage{amsmath}
\usepackage{amsfonts}
\usepackage{amssymb}
\usepackage{graphicx}
\usepackage{geometry}      
\geometry{letterpaper}                  
\usepackage{epstopdf}
\usepackage{fancyhdr} % Paquete para encabezados y pies de pag
\usepackage{listings}
\usepackage{color}
\usepackage{placeins}
\usepackage{csquotes}
\usepackage{textcomp}
\usepackage{gensymb} 

\pagestyle{fancy} 
\chead{\bfseries Fase de formulación} 
\rhead{Grupo C1B. Subgrupo 2}
\lfoot{6 de Mayo de 2023}
\cfoot{Universidad Industrial de Santander} 
\rfoot{\thepage} 

\voffset = -0.25in 
\textheight = 8.0in 
\textwidth = 6.5in
\oddsidemargin = 0.in
\headheight = 20pt 
\headwidth = 6.5in
\renewcommand{\headrulewidth}{0.5pt}
\renewcommand{\footrulewidth}{0,5pt}
\begin{document}
\begin{titlepage}
    \begin{center}
        \includegraphics[width=0.4\textwidth]{../general-images/uis-logo.png}
        
        \vspace{0.5cm}
        \LARGE
        \textbf{Estudio del comportamiento del circuito RLC y el fenómeno de resonancia}
        
        \vspace{0.5cm}
        \large
        Informe 4
        
        \vfill
        
        \textbf{Daniel Esteban Vargas Reyes.} Estudiante - Geología\\
        \textbf{Nicolás Andrés Ramírez Calderón.} Estudiante - Ingeniería de Sistemas\\ 
        \textbf{Rodolfo Valentín Muñoz Vega.} Estudiante - Ingeniería Química\\

        \vspace{1.0cm}
        Presentado a la docente:
        
        \textbf{Zayda Paola Reyes Quijano}
        
        \vfill
        
        Escuela de Física - Física III\\
        Universidad Industrial de Santander\\
        Bucaramanga, Santander, Colombia\\
        20 de Mayo de 2023        
    \end{center}
\end{titlepage}

\tableofcontents

\newpage

\section{Formulación del problema}
En tiempos modernos es común vernos rodeados de una amplia variedad de dispositivos tecnológicos que utilizamos como herramientas para nuestro desarrollo personal y colectivo. Algunas de estas tecnologías funcionan como emisores/receptores de diferentes tipos de ondas, algunas mecánicas (Como
parlantes o micrófonos para sonido) y otras
electromagéticas (Como dispositivos de radiofrecuencias, WiFi, Bluetooth, etc).\par
\bigskip
Nuestro cuerpo no es ajeno a estos fenónemos físicos, ya que también interactuamos con los diferentes tipos de ondas; las cuerdas vocales y el oído tienen un rol en el sonido, mientras que algunas células de nuestra piel se ven claramente afectadas por las ondas electromagnéticas. \cite{IDEAM}\par
\bigskip
En la actualidad, entendemos que las ondas pueden provocar afectaciones en nuestra salud.  Por ejemplo, para ciertas frecuencias somos más sensibles a un alto nivel de intensidad de las ondas sonoras \cite{serway_jewett_2017}, al punto de que puede ser perjudicial para el oído someterse a ondas de
muy elevados decibelios. Por otra parte, las ondas electromagnéticas también pueden provocar efectos negativos, la radiación ultravioleta del sol o los rayos X pueden producir cáncer \cite{Wall2006-pi}.\par
\bigskip
De modo que, es pertinente preguntarnos: ¿Qué características tienen las ondas que pueden producir efectos nocivos en la salud?¿Qué dispotivos comunes pueden producir estas ondas y cuales no?¿Cuales serían las recomendaciones en relación con la exposición a dicho tipo de ondas?\par
\bigskip
Una hipótesis en esta situación podría sugerir que los dispositivos de uso doméstico no emiten ondas nocivas para la salud, ya que no existen impetuosas advertencias sobre el uso de estos dispositivos por parte de los fabricantes o una opinión consensuada por especialistas en salud. Mientras que
quizás algunos aparatos industriales o de uso médico podrían tener algún tipo de afectación en las células de nuestro cuerpo.
\section{Justificación}
\section{Objetivos}
\section{Diseño metodológico}
El desarrollo del proyecto de investigación se puede secuenciar en tres fases principales:
\subsection{Primera fase}
En esta fase se busca ahondar en el entendimiento de los diferentes tipos de ondas y sus efectos en la salud, para ello se requiere de una profundización en los conceptos de ondas de sonido y ondas electromagnéticas.\par
\bigskip
Para esta fase también se pretende realizar una minuciosa consulta sobre el desarrollo académico en los efectos de las ondas en la salud, esto requiere indagar en las producciones academicas realizadas para entender la opinión general de los expertos en el tema.

\subsection{Segunda fase}
Esta fase pretende comprender el funcionamiento de los dispositivos tecnológicos de uso diario, de esta forma se pueden clasificar aquellos que hacen uso de ondas (y también que tipos de ondas utilizan) y descartar los que no hacen uso de ellas.\par
\bigskip
Posteriormente se hará necesario revisar las especificaciones técnicas de los dispositivos previamente mencionados, de esta manera podremos comprender las características de las ondas que son recibidas o emitidas por estos aparatos.
\subsection{Tercera fase}
Para comprobar la información obtenida en la fase anterior se hace necesario medir experimentalmente algunas de las caracterísiticas de las ondas que utilizan los dispositivos domésticos.\par
\bigskip
En esta fase se utilizarán instrumentos de medición que permitan percibir algunas propiedades de las ondas, tales como su amplitud, frecuencia, etc. La información será registrada para su posterior análisis.
\subsection{Cuarta fase}
Por último, se tomará la información obtenida en las fases anteriores para concluir si las ondas empleadas por los dispositivos de uso doméstico representan un riesgo para la salud y en caso de que no lo representen, bajo que circunstancias podrían hacerlo.
\section{Alcances}
\section{Bibliografía} 
\bibliographystyle{unsrt}
\bibliography{references}
\end{document}
